\section{Conclusions and the remaining accuracy conditions}
\label{sec:conclusions}
The implementation provides a multi-step construction: qualified source
generation, canonical first-step proposals, bounded whole-event H
removal/S veto, randomized accepted LHE collection and scale-preserving
continuation. The current sampled mode retains the native history,
Born integration diagram and colour flow before sign classification.
The S event's actual branching identifies its native emission component;
its original integration owner does not constrain the shower radiator.
Diagram and oriented-flow indices are transported under certified
identical-Born-leg permutations, with the corresponding statistical
normalization applied to all densities.

Ordinary matching keeps the generated bogus-Sudakov multiplier:
its transferred numerator is $[P_a(D_{G,a}-R_a)]_+$. Sampled Delta
matching uses the undamped completed numerator $u_a=[D_{G,a}-R_a]_+$
and rejects eligible S events with probability
$\min(u_a/K_a,1)$. The original negative H measure is still
$\Delta_a u_a$, and the residual is $\Delta_a[u_a-K_a]_+$.
This consistent pair removes the need for numerical source-Delta
evaluation at S emissions. It does not replace the G completion by a
scalar shower kernel or divide Delta out of the nominal event weight.
The global-sum compatibility mode retains its original source-damped
ratio and different residual; both prescriptions are explicitly labelled
in the accepted files and reports.

The source G completion, ordinary fixed low-scale factor, physical
shower Sudakov and Delta source factor remain distinct. Hard and MC
mass conventions are recorded separately, and the written first-emission
endpoint carries its own current shower bounds. The finite transfer is
governed by its actual signed S emission density. On the sampled Delta
component space it is
$C_{0,a}[\overline\Delta_a-E_{g,a}^{\Delta}/K_a]$.
Retaining components resolves the nonlinear mixed-history sign problem;
it does not impose equality of the source and physical shower measures.
Negative S coefficients, insufficient support, G/kernel differences,
source-scale mixtures and different continuation laws remain relevant.

The completed ten-million-source-event campaign used the global-sum
prescription. Its count, weight, covariance and kinematic checks passed,
and it exhibits sizeable finite differences for some processes. All
703 figure assets are retained, with campaign labels that prevent their
interpretation as sampled-mode results. Focused native closure, identical
gluon, support-boundary and small-source checks exercise the new
implementation, including genuine negative-H removal and S veto. They
are implementation evidence rather than a new precision production or
a measurement of the sampled prescription's finite distribution defect.

An exact endpoint cancellation would require the actual positive or
signed eligible S density in the same component space, adequate common
support and the corresponding source negative measure. Equality after
further showering additionally requires equality of the transferred
full states under the continuation law. A perturbative accuracy argument
needs the leading native-law/kernel correspondence; its expected order
cannot establish that correspondence or bound logarithmically enhanced
finite effects. The retained source labels, explicit mechanism controls,
portable profiles and complete figure archive make those conditions
stated and testable, without normalizing away the observed differences.
